% Shared reference sheet, included by all three editions.  Every row names a
% term or symbol, states what it means in one line, and points at the place
% where it is properly defined; that pointer is a live link, and in the
% companion editions it leads into the results PDF through xr-hyper.
%
% A row's first cell carries the anchor planted by the \gkey macro: writing
% \term with the same key in the body turns that occurrence into a link to this
% table.  Keep the rows in alphabetical order, and keep every \gkey key matched
% by exactly one \dfn in the body.
%
% A row whose pointer is a companion-owned label must be gated with
%   \ifdefined\SCAFFOLDResultsView\else ... \fi
% or the results edition fails on an undefined reference.  A gated row must not
% be linked from a results module; tools/check_manuscript.pl checks both.

\phantomsection
\section*{Terminology and notation}
\addcontentsline{toc}{section}{Terminology and notation}

Each entry below is defined once in the body, at the place named in the last
column; the definition there is the authoritative one, and what is recorded
here is only enough to read on with.  Terms set in blue elsewhere in the text
link to their row in the first table.

\subsection*{Terminology}

\begingroup
\setlength{\LTleft}{0pt}\setlength{\LTright}{0pt}
\renewcommand{\arraystretch}{1.25}
\begin{longtable}{@{}p{0.24\textwidth}p{0.505\textwidth}p{0.165\textwidth}@{}}
\hline
\textbf{Term} & \textbf{Meaning} & \textbf{Defined in}\\
\hline
\endhead
\hline
\endfoot

\gkey{goldbach-count}{Goldbach count} &
The number $\gc(n)$ of Goldbach pairs for $n$; equivalently, by
Lemma~\ref{reflect:lem}, the size of $\Prset\cap(n-\Prset)$. &
Def.~\ref{pair:def}\\

\gkey{goldbach-number}{Goldbach number} &
An integer admitting at least one Goldbach pair.  Conjecture~\ref{goldbach:conj}
says every even $n\geq4$ is one. &
Def.~\ref{pair:def}\\

\gkey{goldbach-pair}{Goldbach pair} &
An ordered pair $(p,q)$ of primes with $p+q=n$.  Ordered, so $(3,5)$ and
$(5,3)$ are two pairs for $8$. &
Def.~\ref{pair:def}\\

\ifdefined\SCAFFOLDResultsView\else
\gkeyx{hardy-littlewood}{Hardy--Littlewood} &
The conjectural asymptotic $\gc(n)\sim2\Pi_{2}\,n/(\log n)^{2}$ times a
singular series depending on the odd prime factors of $n$.  Used only as a
heuristic; nothing in this project depends on it. &
\S\ref{sec:density}\\
\fi

\gkeyx{sieve}{sieve of Eratosthenes} &
The linear-time marking of composites in $[0,N]$ used by every computation
here.  Standard; stated nowhere in the body. &
\S\ref{sec:verify}\\

\ifdefined\SCAFFOLDResultsView\else
\gkey{sifted-count}{sifted count} &
The count $\gc_{z}(n)$ of $p\leq n$ with neither $p$ nor $n-p$ divisible by a
prime below $z$.  An upper bound for $\gc(n)$, which is the wrong direction. &
Def.~\ref{sifted:def}\\
\fi

\end{longtable}
\endgroup

\subsection*{Notation}

\begingroup
\setlength{\LTleft}{0pt}\setlength{\LTright}{0pt}
\renewcommand{\arraystretch}{1.25}
\begin{longtable}{@{}p{0.24\textwidth}p{0.505\textwidth}p{0.165\textwidth}@{}}
\hline
\textbf{Symbol} & \textbf{Meaning} & \textbf{Defined in}\\
\hline
\endhead
\hline
\endfoot

$\Prset$ & The set of primes. & \S\ref{sec:intro}\\

$n-\Prset$ & $\{n-p:p\in\Prset,\ p\leq n\}$, the reflection of the primes in
$n/2$. & Lem.~\ref{reflect:lem}\\

$\gc(n)$ & The Goldbach count of $n$. & Def.~\ref{pair:def}\\

\ifdefined\SCAFFOLDResultsView\else
$\gc_{z}(n)$ & The sifted count of $n$ at level $z$. & Def.~\ref{sifted:def}\\
\fi

\end{longtable}
\endgroup
